\subsection{Feature-wise Linear Modulation (FiLM)}
\label{subsec:film}

The second model we train also uses convolutional and linear layers to predict the wheel velocities given a segmented input frame, but it relies on the FiLM \cite{Perez2017FiLMVR} conditioning mechanism for the high-level commands.
Given a feature map $x_{\text{in}}$ a FiLM layer learns a shift coefficient $\beta$ and scale coefficient $\gamma$ that augment the input $x_{\text{out}} = \gamma \odot x_{\text{in}} + \beta$.
A convolutional layer is augmented at the channel dimension, each channel of the output feature map is scaled and shifted with a separate learnable shift and scale coefficient. A linear layer only needs one shift and scale pair.
We implement a shared encoder that takes as input the one-hot encoded intent and projects to a 64-dimensional vector. Then for each convolutional and linear layer that is augmented with FiLM we have a seperate projection that outputs the correct dimensionality.
For all FiLM layers we initialize $\gamma = 1$ and $\beta = 0$.
A detailed overview of the FiLM-conditioned neural network can be found in Figure \ref{fig:FiLM}.

